% =============================================================================
% Chapter 19: OLAP και Data Warehousing - Λυμένες Ασκήσεις
% Data Mining Study Guide - NTUA
% =============================================================================

\chapter{OLAP και Data Warehousing - Λυμένες Ασκήσεις}
\label{ch:olap-dw-exercises}

Αυτό το κεφάλαιο περιέχει λυμένες ασκήσεις για OLAP operations, star schema design,
και data warehouse queries.

% -----------------------------------------------------------------------------
% Exercise 8.1
% -----------------------------------------------------------------------------
\section{Άσκηση 8.1: OLAP Operations Ταξινόμηση}

\begin{formulabox}[Εκφώνηση]
Ταξινομήστε κάθε λειτουργία σε Roll-up, Drill-down, Slice, Dice, ή Pivot:

\begin{enumerate}
    \item Από πωλήσεις ανά ημέρα σε πωλήσεις ανά μήνα
    \item Εμφάνιση μόνο των πωλήσεων του 2024
    \item Από πωλήσεις ανά χώρα σε πωλήσεις ανά πόλη
    \item Πωλήσεις για συγκεκριμένη περιοχή ΚΑΙ κατηγορία προϊόντος
    \item Αλλαγή από Product×Time σε Region×Time
\end{enumerate}
\end{formulabox}

\subsection*{Λύση}

\begin{center}
\begin{tabular}{clp{6cm}}
\toprule
\textbf{\#} & \textbf{Operation} & \textbf{Εξήγηση} \\
\midrule
1 & \textbf{Roll-up} & Aggregation σε υψηλότερο επίπεδο ιεραρχίας (ημέρα → μήνας) \\
2 & \textbf{Slice} & Επιλογή μιας τιμής σε μία διάσταση (year = 2024) \\
3 & \textbf{Drill-down} & Μετάβαση σε χαμηλότερο επίπεδο λεπτομέρειας \\
4 & \textbf{Dice} & Επιλογή εύρους τιμών σε πολλές διαστάσεις \\
5 & \textbf{Pivot} & Αναδιάταξη διαστάσεων/αλλαγή οπτικής γωνίας \\
\bottomrule
\end{tabular}
\end{center}

\begin{center}
\begin{tikzpicture}[scale=0.7]
    % Hierarchy visualization
    \node[draw, fill=box-blue, minimum width=3cm] at (0,3) {Year};
    \node[draw, fill=box-blue, minimum width=3cm] at (0,2) {Quarter};
    \node[draw, fill=box-blue, minimum width=3cm] at (0,1) {Month};
    \node[draw, fill=box-blue, minimum width=3cm] at (0,0) {Day};

    \draw[arrow, thick, success] (1.8,0.5) -- (1.8,2.5) node[midway, right] {Roll-up};
    \draw[arrow, thick, accent] (-1.8,2.5) -- (-1.8,0.5) node[midway, left] {Drill-down};
\end{tikzpicture}
\end{center}

% -----------------------------------------------------------------------------
% Exercise 8.2
% -----------------------------------------------------------------------------
\section{Άσκηση 8.2: Roll-up και Drill-down Sequence}

\begin{formulabox}[Εκφώνηση]
Δίνεται cube με διαστάσεις: Time (Day → Month → Quarter → Year),
Product (Item → Category → Department), Location (City → Region → Country).

Ξεκινώντας από (Day, Item, City), περιγράψτε τη σειρά operations για να φτάσετε σε
(Year, Department, Country).
\end{formulabox}

\subsection*{Λύση}

\textbf{Απαιτούμενα Roll-ups:}

\begin{enumerate}
    \item \textbf{Time}: Day → Month → Quarter → Year (3 roll-ups)
    \item \textbf{Product}: Item → Category → Department (2 roll-ups)
    \item \textbf{Location}: City → Region → Country (2 roll-ups)
\end{enumerate}

\textbf{Σύνολο:} 7 consecutive roll-up operations

\begin{center}
\begin{tikzpicture}[node distance=0.8cm, font=\small]
    \node[draw, fill=box-green, rounded corners] (start) {(Day, Item, City)};
    \node[draw, fill=box-blue, rounded corners, below=of start] (s1) {(Month, Item, City)};
    \node[draw, fill=box-blue, rounded corners, below=of s1] (s2) {(Quarter, Category, Region)};
    \node[draw, fill=box-green, rounded corners, below=of s2] (end) {(Year, Department, Country)};

    \draw[arrow] (start) -- (s1) node[midway, right] {roll-up Time};
    \draw[arrow] (s1) -- (s2) node[midway, right] {roll-up all};
    \draw[arrow] (s2) -- (end) node[midway, right] {roll-up all};
\end{tikzpicture}
\end{center}

\begin{infobox}[Optimization]
Στην πράξη, τα roll-ups μπορούν να γίνουν παράλληλα σε διαφορετικές διαστάσεις
αν υπάρχουν pre-computed aggregates.
\end{infobox}

% -----------------------------------------------------------------------------
% Exercise 8.3
% -----------------------------------------------------------------------------
\section{Άσκηση 8.3: Slice vs Dice}

\begin{formulabox}[Εκφώνηση]
Για cube πωλήσεων με διαστάσεις (Product, Location, Time), εκφράστε σε OLAP
operations:

\begin{enumerate}
    \item Πωλήσεις για τον Ιανουάριο 2024
    \item Πωλήσεις για Electronics σε Athens τον Ιανουάριο 2024
    \item Πωλήσεις για όλα τα προϊόντα, στην Ελλάδα, για Q1 2024
\end{enumerate}
\end{formulabox}

\subsection*{Λύση}

\textbf{(1) Slice}

$$\text{SLICE}(Cube, Time = \text{``Jan 2024''})$$

Μειώνει τον cube κατά μία διάσταση.

\textbf{(2) Dice}

$$\text{DICE}(Cube, Product = \text{``Electronics''}, Location = \text{``Athens''}, Time = \text{``Jan 2024''})$$

Επιλογή σε πολλαπλές διαστάσεις ταυτόχρονα.

\textbf{(3) Slice + Roll-up (ή Dice)}

\begin{enumerate}
    \item Roll-up Location: City → Country (Ελλάδα)
    \item Roll-up Time: Month → Quarter (Q1)
    \item Slice: Country = ``Greece'', Quarter = ``Q1 2024''
\end{enumerate}

% -----------------------------------------------------------------------------
% Exercise 8.4
% -----------------------------------------------------------------------------
\section{Άσκηση 8.4: Cube Calculations}

\begin{formulabox}[Εκφώνηση]
Ένας cube έχει 3 διαστάσεις: Product (100 values), Location (50 values),
Time (365 days).

\begin{enumerate}
    \item Πόσα cells έχει ο base cube;
    \item Πόσα συνολικά cells έχει ο πλήρης CUBE (με όλα τα aggregations);
    \item Αν κάθε cell χρειάζεται 8 bytes, πόση μνήμη απαιτείται;
\end{enumerate}
\end{formulabox}

\subsection*{Λύση}

\textbf{(1) Base Cube Cells:}

$$|Product| \times |Location| \times |Time| = 100 \times 50 \times 365 = 1,825,000$$

\textbf{(2) Full CUBE:}

Για κάθε διάσταση, μπορεί να έχουμε τις τιμές της ή το ALL (aggregation):

$$(|Product| + 1) \times (|Location| + 1) \times (|Time| + 1)$$
$$= 101 \times 51 \times 366 = 1,884,666$$

Εναλλακτικά (για cuboids):
$$\text{Cuboids} = 2^3 = 8$$

\begin{center}
\begin{tabular}{lc}
\toprule
\textbf{Cuboid} & \textbf{Cells} \\
\midrule
(P, L, T) & $100 \times 50 \times 365 = 1,825,000$ \\
(P, L, *) & $100 \times 50 = 5,000$ \\
(P, *, T) & $100 \times 365 = 36,500$ \\
(*, L, T) & $50 \times 365 = 18,250$ \\
(P, *, *) & $100$ \\
(*, L, *) & $50$ \\
(*, *, T) & $365$ \\
(*, *, *) & $1$ \\
\midrule
\textbf{Total} & \textbf{1,885,266} \\
\bottomrule
\end{tabular}
\end{center}

\textbf{(3) Memory:}

$$1,885,266 \times 8 \text{ bytes} = 15,082,128 \text{ bytes} \approx 14.4 \text{ MB}$$

% -----------------------------------------------------------------------------
% Exercise 8.5
% -----------------------------------------------------------------------------
\section{Άσκηση 8.5: Star Schema Design}

\begin{formulabox}[Εκφώνηση]
Σχεδιάστε star schema για ηλεκτρονικό κατάστημα με τα εξής requirements:
\begin{itemize}
    \item Παρακολούθηση πωλήσεων (ποσότητα, τιμή, έκπτωση)
    \item Διαστάσεις: Προϊόν, Πελάτης, Ημερομηνία, Κατάστημα
    \item Ιεραρχία προϊόντων: Product → Subcategory → Category → Department
\end{itemize}
\end{formulabox}

\subsection*{Λύση}

\begin{center}
\begin{tikzpicture}[
    fact/.style={rectangle, draw=accent, fill=box-red, minimum width=3cm, minimum height=2cm, font=\small, align=center},
    dim/.style={rectangle, draw=secondary, fill=box-blue, minimum width=2.5cm, minimum height=1.8cm, font=\small, align=center}]

    % Fact table
    \node[fact] (fact) {
        \textbf{Sales\_Fact}\\
        \underline{product\_sk}\\
        \underline{customer\_sk}\\
        \underline{date\_sk}\\
        \underline{store\_sk}\\
        quantity\\
        unit\_price\\
        discount
    };

    % Dimension tables
    \node[dim, above left=1cm and 2cm of fact] (product) {
        \textbf{Dim\_Product}\\
        \underline{product\_sk}\\
        product\_id\\
        name\\
        subcategory\\
        category\\
        department
    };

    \node[dim, above right=1cm and 2cm of fact] (customer) {
        \textbf{Dim\_Customer}\\
        \underline{customer\_sk}\\
        customer\_id\\
        name\\
        city\\
        country
    };

    \node[dim, below left=1cm and 2cm of fact] (date) {
        \textbf{Dim\_Date}\\
        \underline{date\_sk}\\
        date\\
        day\_of\_week\\
        month\\
        quarter\\
        year
    };

    \node[dim, below right=1cm and 2cm of fact] (store) {
        \textbf{Dim\_Store}\\
        \underline{store\_sk}\\
        store\_id\\
        name\\
        city\\
        region
    };

    % Connections
    \draw[thick] (product) -- (fact);
    \draw[thick] (customer) -- (fact);
    \draw[thick] (date) -- (fact);
    \draw[thick] (store) -- (fact);
\end{tikzpicture}
\end{center}

\begin{infobox}[Surrogate Keys (\_sk)]
Χρησιμοποιούμε surrogate keys (τεχνητά κλειδιά) αντί για natural keys για:
\begin{itemize}
    \item Καλύτερη απόδοση (μικρότερα integers)
    \item Διαχείριση slowly changing dimensions
    \item Ανεξαρτησία από source systems
\end{itemize}
\end{infobox}

% -----------------------------------------------------------------------------
% Exercise 8.6
% -----------------------------------------------------------------------------
\section{Άσκηση 8.6: OLAP Queries σε SQL}

\begin{formulabox}[Εκφώνηση]
Χρησιμοποιώντας το star schema της προηγούμενης άσκησης, γράψτε SQL queries για:

\begin{enumerate}
    \item Συνολικές πωλήσεις ανά κατηγορία προϊόντος
    \item Top 5 πελάτες με βάση το συνολικό ποσό αγορών
    \item Πωλήσεις ανά μήνα για το 2024
\end{enumerate}
\end{formulabox}

\subsection*{Λύση}

\textbf{(1) Πωλήσεις ανά κατηγορία (Roll-up):}

\begin{lstlisting}[language=SQL]
SELECT
    p.category,
    SUM(f.quantity * f.unit_price * (1 - f.discount)) AS total_revenue
FROM Sales_Fact f
JOIN Dim_Product p ON f.product_sk = p.product_sk
GROUP BY p.category
ORDER BY total_revenue DESC;
\end{lstlisting}

\textbf{(2) Top 5 πελάτες:}

\begin{lstlisting}[language=SQL]
SELECT
    c.name,
    SUM(f.quantity * f.unit_price * (1 - f.discount)) AS total_spent
FROM Sales_Fact f
JOIN Dim_Customer c ON f.customer_sk = c.customer_sk
GROUP BY c.customer_sk, c.name
ORDER BY total_spent DESC
LIMIT 5;
\end{lstlisting}

\textbf{(3) Πωλήσεις ανά μήνα 2024 (Slice + Roll-up):}

\begin{lstlisting}[language=SQL]
SELECT
    d.year,
    d.month,
    SUM(f.quantity * f.unit_price * (1 - f.discount)) AS monthly_revenue
FROM Sales_Fact f
JOIN Dim_Date d ON f.date_sk = d.date_sk
WHERE d.year = 2024
GROUP BY d.year, d.month
ORDER BY d.month;
\end{lstlisting}

\begin{warningbox}[Performance]
Στα data warehouses, οι dimension tables είναι συνήθως denormalized για αποφυγή
πολλαπλών JOINs. Αυτό αυξάνει το μέγεθος αλλά βελτιώνει την απόδοση των queries.
\end{warningbox}
